% from ega4-2-fr.tex
\oldpage[IV-2]{210}

la prop. 9), donc
$\widehat B\otimes_BK'=\widehat{B'}\otimes_{B'}K'$ par associativité du produit
tensoriel; mais $B'$ est une $A$-algèbre intègre finie, donc $\widehat{B'}$ est
un anneau réduit en vertu de $a)$, et comme $\widehat{B'}$ est un $B'$-module
plat, les éléments $\neq0$ de $B'$ sont $\widehat{B'}$-réguliers, ce qui
entraîne que $\widehat{B'}\otimes_{B'}K'$ s'identifie à un sous-anneau de
l'anneau total des fractions de $\widehat{B'}$, et a fortiori est réduit.

Pour voir que $a')$ entraîne $a'')$, considérons un point quelconque $x$ de
$X=\operatorname{Spec}(A)$, et soit $Y$ le sous-schéma fermé intègre de $X$
ayant $\overline{\{x\}}$ pour espace sous-jacent; on a
$Y=\operatorname{Spec}(B)$, où $B$ est un anneau quotient intègre de $A$, et
$\operatorname{Spec}(\widehat B)=Y\times_X\operatorname{Spec}(\widehat A)$,
de sorte que la fibre formelle de $A$ au point $x$ est la même que celle de
$B$, ou encore est égale à $\operatorname{Spec}(\widehat B\otimes_BK)$. Comme
les anneaux locaux de $\operatorname{Spec}(\widehat B\otimes_BK)$ sont ceux de
$\operatorname{Spec}(\widehat B)$ aux points de la fibre de $x$, l'hypothèse
entraîne que $\widehat B\otimes_BK$ est réduit, et les conditions de $a')$
entraînent donc que $\widehat B\otimes_BK$ est une $K$-algèbre séparable
(4.6.1).

La condition $a'')$ entraîne $a)$, car il résulte de (7.3.15) que les fibres
formelles de $C$ sont alors géométriquement réduites, et si $C$ est réduit, il
en est donc de même de $\widehat C$ (cas particulier de (7.3.17)).

Le fait que $a')$ implique $b)$ est un cas particulier de (0, 23.1.7). Reste
donc à prouver que $b)$ entraîne $a)$. Notons que si $C$ est une $A$-algèbre
finie, pour tout idéal premier $\mathfrak q$ de $C$, l'image réciproque
$\mathfrak p$ de $\mathfrak q$ dans $A$ est un idéal premier et
$C/\mathfrak q$ est une $(A/\mathfrak p)$-algèbre finie, donc l'hypothèse $b)$
entraîne que tout anneau quotient intègre de $C$ est un anneau japonais. On est
ainsi ramené à prouver que sous l'hypothèse $b)$, le complété de tout quotient
intègre de $A$ est réduit. Raisonnons par récurrence sur $n=\dim(A)$,
l'assertion étant triviale pour $n=0$; en remplaçant $A$ par les quotients de
$A$ par ses idéaux premiers minimaux $\mathfrak p_i$, on peut se borner au cas
où $A$ est intègre (tout quotient intègre de $A$ étant quotient de l'un des
$A/\mathfrak p_i$). Pour tout idéal premier $\mathfrak p\neq0$, l'hypothèse de
récurrence montre déjà que le complété de $A/\mathfrak p$ est réduit, et il
suffit donc de prouver que $\widehat A$ est réduit. En outre, la clôture
intégrale $A'$ de $A$ est par hypothèse un $A$-module de type fini, donc un
anneau semi-local noethérien et $\widehat A$ s'identifie à un sous-anneau de
$\widehat{A'}$ ((0$_{\mathrm I}$, 7.3.3) et Bourbaki, \emph{Alg. comm.}, chap.
IV, § 2, n\textsuperscript{o} 5, cor. 3 de la prop. 9); il suffira donc de
prouver que $\widehat{A'}$ est réduit; on a vu ci-dessus que l'hypothèse $b)$
est aussi vérifiée par $A'$, qui par ailleurs est aussi de dimension $n$
(0, 16.1.5); on peut donc se borner au cas où $A$ est intégralement clos. Soit
$t\neq0$ un élément du radical de $A$, et soient $\mathfrak q_j$
($1\leq j\leq n$) les idéaux premiers minimaux parmi ceux contenant $tA$; on a
les propriétés suivantes :
\begin{enumerate}
\item[\emph{(i)}] $t$ est régulier.
\item[\emph{(ii)}] $A/tA$ n'a pas d'idéaux premiers associés immergés.
\item[\emph{(iii)}] Les $A_{\mathfrak q_j}$ sont des anneaux de valuation
discrète.
\item[\emph{(iv)}] Les complétés des $A/\mathfrak q_j$ sont réduits.
\end{enumerate}
En effet, (i) est triviale puisque $A$ est intègre et $t\neq0$. Comme $A$ est
intégralement clos, $A/tA$ vérifie $(\mathrm S_1)$, c'est-à-dire (5.7.7) est
sans idéaux premiers associés immergés (Bourbaki, \emph{Alg. comm.}, chap. VII,
§ 1, n\textsuperscript{o} 4, prop. 8). Toujours parce que $A$ est
intégralement

